\documentclass[12pt]{article}

\usepackage[utf8]{inputenc}
\usepackage[uzbek]{babel}
\usepackage{amsmath, amssymb, amsthm}
\usepackage{geometry}
\usepackage{graphicx}
\usepackage{tikz}
\usepackage{booktabs}
\usepackage{caption}
\usepackage{fontawesome5}
\usepackage{tcolorbox}
\usetikzlibrary{shapes, backgrounds}

\addto\captionsuzbek{\renewcommand{\tablename}{Jadval}}

\geometry{a4paper, margin=1in}

\title{Fazzi to'plamlar}
\author{Shahrizoda Kenjaboyeva}
\date{\today}

\begin{document}
	
	\maketitle
	
	\section{Fazzi to'plam tushunchasi}
	
	$(\mu_A(x), x)$ ko'rinishidagi juftliklar to'plamiga \textbf{A fazzi to'plami} deyiladi.
	Bu yerda $x \in X$ va $\mu_A(x)$-$x$ ning $A$ to'plamga a'zolik funksiyasi.
	
	\[\mu_A(X): X \to [0, 1]\]
	
	Agar $B \subseteq X$ bo'lsa, uning a'zolik funksiyasi quyidagicha aniqlanadi:
	\[\mu_B(x)=
	\begin{cases}
		1, & x\in B \\
		0, & x\notin B
	\end{cases}\]
	Bundan ko'rinadiki oddiy to'plamlar fazzi to'plamlarning xususiy holidir.
	
	\subsection{Asosiy tushunchalar}
	
	\begin{itemize}
		\item \textbf{Tashuvchi:} $\mu_A(x)>0$ bo'lgan barcha $x$ lar $A$ ning tashuvchilari deyiladi.
		\item \textbf{Bir nuqtali fazzi to'plam:} Tashuvchisi faqat bitta elementdan iborat bo'lgan fazzi to'plam: $A=\frac{\mu}{x}$.
		\item \textbf{O'tish nuqtasi:} $\mu_A(x)=0.5$ bo'lgan $x$ nuqta.
	\end{itemize}
	
	\section{Fazzi to'plamlarni ifodalash}
	
	\begin{itemize}
		\item $A=\int \frac{\mu_A(x)}{x} \, dx$
		
		Tashuvchilari chekli elementlardan iborat $A$ fazzi to'plamni quyidagicha belgilash ham mumkin:
		
		\item $A=\frac{\mu_1}{x_1}+\frac{\mu_2}{x_2}+...+\frac{\mu_n}{x_n}$
		$A=\sum_{i=1}^{n}\frac{\mu_i}{x_i}$
		
		\begin{table}[h]
			\centering
			\caption{Fazzi to'plamlarni ifodalashning yana bir usuli:}
			\vspace{2mm}
			\begin{tabular}{c|c|c|c|c}
				\toprule
				$\mu_i$ & $\mu_1$ & $\mu_2$ & ...& $\mu_n$ \\ \midrule
				$x_i$ & $x_1$ & $x_2$ & ... &$x_n$ \\ \bottomrule
			\end{tabular}
		\end{table}
		
	\end{itemize}
	
	\begin{quote}
		\faExclamationTriangle\ \textbf{Eslatma:} Bularning hech biri integral, summa yoki taqsimot jadvali emas!
	\end{quote}


Masalan, $A=\frac{0.2}{3}+\frac{0.2}{4}+\frac{0.5}{6}+\frac{0.7}{7}+\frac{0.3}{8}$

\begin{itemize}
	\item $\sup_{a\in X}\mu_A(x)$ - $A$ fazzi to'plam balandligi;
	\item Balandligi $1$ ga teng fazzi to'plam - normal fazzi to'plam;
	\item $1$ dan kichik - subnormal fazzi to'plam:
	\item Faqat bitta tashuvchisining a'zolik darajasi $1$ ga teng to'plam - unimodal fazzi to'plam;
	\item Subnormalni normalga aylantirish:
	
	\[\mu_A(x):=\frac{\mu_A(x)}{\sup_{x\in X}\mu_A(x)}\]
	
	\item $A$ va $B$ $X$ universal to'plamda a'zolik funksiyalari $\mu_A(x)$ va $\mu_B(x)$ bo'lgan fazzi to'plamlari bo'lsin.
	Agar har bir $x\in X$ element uchun $\mu_A(x)=\mu_B(x)$ bo'lsa, $A$ va $B$ fazzi to'plamlar teng deyiladi. Belgilanishi:

	\[A=B\]
	
	Tenglik o'rniga ekvivalentlik ham qo'llaniladi:
	
	\[A\sim B\]
	
	\item Agar barcha $x\in X$ uchun $\mu_A(x)=0$ bo'lsa, $A$ bo'sh fazzi to'plam deyiladi va $A=\varnothing $ kabi belgilanadi.
	
	\item Agar barcha $x\in X$ uchun $\mu_A(x)\leqslant \mu_B(x)$ bo'lsa, $A$ fazzi to'plam $B$ fazzi to'plamning qismi deyiladi: $A\subseteq B$
	
	\item Agar $\mu_B(x)=1-\mu_A(x)$ bo'lsa, $A$ va $B$ fazzi to'plamlar bir-birini to'ldiradi deyiladi: $B=\bar{A}$
\end{itemize}

\subsection{Fazzi to'plamlar ustida amallar}

\begin{itemize}
	\item $\mu_{A\cup B}(x)=\max{\mu_A(x), \mu_B(x)}$ - $A \cup B$ fazzi to'plamning a'zolik funksiyasi;
	
	\item $\mu_{A\cap B}(x)=\min{\mu_A(x), \mu_B(x)}$ - $A \cap B$ fazzi to'plamning a'zolik funksiyasi;
	
	\item $A\setminus B$ fazzi to'plamning a'zolik funksiyasi:
	\[
	\mu_{A\setminus B}(x)=
	\begin{cases}
		\mu_A(x)-\mu_B(x), & \text{agar }\mu_A(x)\geq \mu_B(x)\text{ bo'lsa} \\
		0, & \text{aks holda}
	\end{cases}
	\] 
\end{itemize}


	
\end{document}\documentclass{article}